\documentclass[12pt,a4paper]{article}
\usepackage{amsmath,amssymb,amsthm}
\usepackage{hyperref}
\usepackage{geometry}
\usepackage{enumitem}
\geometry{margin=2.5cm}
\usepackage{booktabs}

\newtheorem{theorem}{Theorem}[section]
\newtheorem{lemma}[theorem]{Lemma}
\newtheorem{corollary}[theorem]{Corollary}
\newtheorem{proposition}[theorem]{Proposition}
\theoremstyle{definition}
\newtheorem{definition}[theorem]{Definition}
\theoremstyle{remark}
\newtheorem{remark}[theorem]{Remark}

\title{The Gauge Group
$SU(3) \times SU(2) \times U(1)$\\
from Recognition Science:\\
Dimension-Forced Ranks and Structure}

\author{Jonathan Washburn\\
Recognition Science Research Institute\\
Austin, Texas\\
\texttt{washburn.jonathan@gmail.com}}

\date{March 2026}

\begin{document}
\maketitle

\begin{abstract}
We derive the Standard Model gauge group
$SU(3)_C \times SU(2)_L \times U(1)_Y$ from the Recognition
Science (RS) forcing chain~\cite{RS_Axioms}.  The group ranks are
forced by $D = 3$ spatial dimensions: $N_c = 3$ colors from the
face-pair count of the $D$-cube ($F/2 = D = 3$), $SU(2)$ from the
spinor structure of the rotation group $SO(3)$, and $U(1)$ from
the overall recognition phase.  We prove why unitary groups
arise (not orthogonal or symplectic), verify anomaly
cancellation for three generations, and contrast RS with grand
unified theories.  The gauge group is not postulated but derived
from the geometry of the 3-cube and the 8-tick epoch structure.
\end{abstract}

%% ============================================================
\section{Introduction}
\label{sec:intro}
%% ============================================================

The Standard Model gauge group
$SU(3)_C \times SU(2)_L \times U(1)_Y$ was constructed
empirically (Glashow~\cite{Glashow1961},
Weinberg~\cite{Weinberg1967},
Salam~\cite{Salam1968}) but is postulated, not derived.
Several deep questions remain unanswered:
\begin{enumerate}
  \item Why $SU(3)$ for color?  Why three colors, not two or
  four?
  \item Why $SU(2)$ for weak isospin?  Why not $SU(3)$ or
  $U(1)$?
  \item Why $U(1)$ for hypercharge?
  \item Why unitary groups?  Why not $SO(n)$ or $Sp(n)$?
  \item Why this specific direct product structure?
\end{enumerate}

Recognition Science (RS)~\cite{RS_Axioms} provides an alternative:
the gauge group emerges from the $D=3$ cube geometry rather than
from a postulated large group.

Grand unified theories (GUTs) such as
$SU(5)$~\cite{GeorgiGlashow1974},
$SO(10)$~\cite{Georgi1975}, or $E_8$~\cite{Lisi2007} embed
the SM into larger groups but must invoke spontaneous symmetry
breaking at a GUT scale, introducing additional free parameters
and predicting proton decay not yet observed.  RS provides a
bottom-up alternative: the SM gauge group and its representations
emerge from the combinatorial structure of the $D = 3$ recognition
cube without invoking a larger group, symmetry breaking, or extra
dimensions.

%% ============================================================
\section{Recognition Science Framework}
\label{sec:framework}
%% ============================================================

Recognition Science derives all physics from a single primitive,
the \emph{recognition cost functional} $J(x)$, uniquely determined
by three axioms.

\begin{definition}[RS Cost Functional]
For $x > 0$:
\begin{equation}
  J(x) = \frac{1}{2}\!\left(x + x^{-1}\right) - 1.
\end{equation}
\end{definition}

\noindent\textbf{Axiom A1 (Normalization):} $J(1) = 0$.\\
\textbf{Axiom A2 (Recognition Composition Law, RCL):}
$J(xy) + J(x/y) = 2J(x)J(y) + 2J(x) + 2J(y)$ for all $x, y > 0$.\\
\textbf{Axiom A3 (Calibration):} $J''_{\log}(0) = 1$, where
$J_{\log}(t) := J(e^t)$.

\begin{theorem}[Cost Uniqueness]
\label{thm:unique}
The unique continuous solution to \textup{(A1)--(A3)} is
$J(x) = \tfrac{1}{2}(x + x^{-1}) - 1$.
\end{theorem}

\begin{proof}[Proof sketch]
Setting $f(t) = J_{\log}(t) + 1$ and rewriting \textup{(A2)} gives
the d'Alembert functional equation $f(s+t) + f(s-t) = 2f(s)f(t)$.
By the Acz\'el classification theorem~\cite{Aczel1966}, continuous
solutions with $f(0) = 1$ and $f''(0) = 1$ are $f(t) = \cosh t$,
yielding $J(x) = \tfrac{1}{2}(x + x^{-1}) - 1$.
\end{proof}

\noindent The key symmetry is $J(x) = J(1/x)$ (reciprocal symmetry).
The 8-tick epoch ($2^D = 8$ for $D = 3$, forced by gap-45 synchronization)
gives the ledger its complex structure; this is the origin of the unitary
(rather than orthogonal or symplectic) gauge groups.  Particle masses are
$\varphi$-rung assignments: $m = Y_s \cdot E_{\rm coh} \cdot \varphi^r$.

%% ============================================================
\section{The $D$-Cube Combinatorics}
\label{sec:cube}
%% ============================================================

\begin{definition}[$D$-Cube Invariants]
\label{def:cube}
A $D$-dimensional hypercube has the following combinatorial
invariants:
\begin{align}
  V &= 2^D & &\text{(vertices)}, \\
  E &= D \cdot 2^{D-1} & &\text{(edges)}, \\
  F &= 2D & &\text{(faces, i.e., $(D-1)$-faces)}, \\
  F_2 &= \binom{D}{2} & &\text{(2-dimensional faces)}.
\end{align}
For $D = 3$: $V = 8$, $E = 12$, $F = 6$, $F_2 = 3$.
\end{definition}

\begin{definition}[Face Pairs]
\label{def:face-pairs}
A \emph{face pair} of the $D$-cube is a pair of opposite
$(D-1)$-faces, or equivalently, a choice of a coordinate axis.
The number of face pairs equals $D$.  A \emph{2-plane pair} is
a choice of 2 coordinate axes from $D$, specifying a 2D
coordinate plane.  The count of 2-plane pairs is
$\binom{D}{2} = D(D-1)/2$.
\end{definition}

\begin{lemma}[Face-Pair Counts for $D = 3$]
\label{lem:d3}
For $D = 3$:
\begin{enumerate}[label=(\roman*)]
  \item Face pairs: $3$ (the $xy$, $xz$, $yz$ opposite-face
  pairs).
  \item 2-plane pairs: $\binom{3}{2} = 3$ (the coordinate
  planes).
\end{enumerate}
Both counts equal $3$, which is specific to $D = 3$.
\end{lemma}

%% ============================================================
\section{Gauge Ranks from $D = 3$}
\label{sec:ranks}
%% ============================================================

\begin{proposition}[Color: $N_c = 3$ from Face Pairs --- Structural Identification]
\label{thm:nc}
The RS identification is that the number of quark colors equals the number of
face pairs of the 3-cube:
\begin{equation}
  N_c = \frac{F}{2} = \frac{2D}{2} = D = 3.
  \label{eq:nc}
\end{equation}
For $D = 3$ specifically, this also equals $\binom{D}{2} = 3$,
the number of independent 2-plane pairs; the two counts
coincide only at $D = 3$ (see Lemma~\ref{lem:d3}).
\end{proposition}

\begin{proof}[Structural argument]
Each face pair of the 3-cube defines an independent
``recognition channel'' on the ledger.  The RS identification is that a quark's
color state specifies which face-pair channel carries its recognition events.
Since the 3-cube has exactly 3 face pairs, there are exactly 3 independent
color channels.  For $D = 2$, $F/2 = 2$ (only 2 colors); for $D = 4$,
$F/2 = 4$ (4 colors).  Since $D = 3$ is forced by the 8-tick epoch and
gap-45 synchronization, $N_c = 3$ is uniquely determined under this
identification.  The gap between this combinatorial argument and a full
derivation of $SU(3)$ as the gauge group (including Lie algebra structure,
coupling strength, and confinement) is an identified open problem.
\end{proof}

\begin{proposition}[Weak Isospin: $SU(2)$ from Doublet Symmetry --- Structural Identification]
\label{thm:su2}
The RS identification is that the weak $SU(2)$ gauge symmetry arises from the
doublet structure of left-handed fermion recognition events on the ledger.
\end{proposition}

\begin{proof}[Structural argument]
In the RS ledger, each 8-tick epoch partitions elementary
excitations into 4 forward ticks and 4 backward ticks; left-handed
fermions are structurally assigned to the forward half.  The minimal
symmetry group preserving the inner product and determinant on a
2-dimensional complex space (the ``up/down'' doublet of left-handed
quark or lepton partners) is $SU(2)$.  The binary split of the 8
tick phases into two groups of 4 motivates 2 isospin states per
doublet, forcing the gauge rank to be 2.  The resulting gauge group is
identified with $SU(2)_L$.  The identification of the RS 4+4 tick split
with left-handed chirality is the key structural postulate; a derivation of
left-handedness of weak interactions from the RS Hamiltonian is an open problem
(see Remark~\ref{rem:su2-distinction}).
\end{proof}

\begin{remark}[Weak $SU(2)$ versus spin $SU(2)$]
\label{rem:su2-distinction}
The weak $SU(2)$ gauge group derived above is an \emph{internal}
symmetry (isospin) that acts differently from the spatial rotation
$SU(2)$.  Spatial rotations form the covering group
$SU(2) \to SO(3)$, which acts on spin-$1/2$ states.  The weak
$SU(2)$ acts on \emph{isospin} doublets (e.g., $(u, d)_L$ or
$(\nu_e, e)_L$), not on spatial spin.  These are distinct groups,
though both are $SU(2)$ algebraically.  The RS derivation above
gives the isospin doublet structure from the 8-tick phase splitting;
the connection to spatial rotations (which are also constrained by the
$D = 3$ geometry) is a structural coincidence, not a confusion.
A more rigorous derivation of the left-handedness of weak interactions
from the RS ledger asymmetry is an identified open problem.
\end{remark}

\begin{proposition}[Hypercharge: $U(1)$ from Overall Phase --- Structural Identification]
\label{thm:u1}
The RS identification is that the $U(1)_Y$ hypercharge symmetry arises from the
overall recognition phase of ledger events.
\end{proposition}

\begin{proof}[Structural argument]
Each recognition event on the ledger carries a complex phase
$e^{i\theta}$ (from the 8-tick phase structure).  The $J$-cost
$J(x) = J(|x|)$ is invariant under global phase rotations
$x \to e^{i\theta} x$.  This global $U(1)$ phase symmetry is
structurally identified with weak hypercharge, and gauging it
(promoting $\theta$ to a local function $\theta(\mathbf{n})$) gives
the hypercharge field $B_\mu$.  The step from ``$J$-cost is phase-invariant''
to ``this phase must be gauged as $U(1)_Y$'' requires identifying the
RS phase with the weak hypercharge quantum number, which assigns
specific values (e.g., $Y = 1/6$ to $Q_L$, $Y = 1$ to $e_R^c$) from
the gap function $\mathrm{gap}(Z) = Z + Z^2 + Z^4$.  The derivation of
these specific hypercharge assignments from the RS mass formula is an
identified open problem.
\end{proof}

\begin{corollary}[Standard Model Gauge Group]
\label{cor:sm}
The full gauge group is:
\begin{equation}
  G_{\mathrm{SM}} = SU(3)_C \times SU(2)_L \times U(1)_Y,
\end{equation}
with ranks $(3, 2, 1)$ forced by $D = 3$.
\end{corollary}

%% ============================================================
\section{Why Unitary Groups}
\label{sec:unitary}
%% ============================================================

\begin{proposition}[Unitarity from Complex Ledger Structure --- Structural Argument]
\label{thm:unitary}
The RS structural argument is that gauge groups are unitary ($SU(n)$ and $U(1)$),
not orthogonal ($SO(n)$) or symplectic ($Sp(n)$), because the
ledger is fundamentally complex.
\end{proposition}

\begin{proof}
The RS ledger encodes recognition events as complex ratios
$x \in \mathbb{C}^*$, arising from the 8-tick phase structure:
the primitive 8th root of unity $\omega = e^{i\pi/4}$
generates the cyclic group $C_8$, which acts on the complex
plane.  Transformations that preserve the recognition structure
must preserve the complex inner product on the space of ledger
states (i.e., be unitary).

\begin{itemize}
  \item $SO(n)$ preserves a real inner product.  But the
  ledger is complex (the 8-tick forces $\mathbb{C}$, not
  $\mathbb{R}$; see the complex number derivation in a
  companion paper), so $SO(n)$ is too restrictive.
  \item $Sp(n)$ preserves an antisymmetric bilinear form.
  The $J$-cost is symmetric ($J(x) = J(x^{-1})$), not
  antisymmetric, so $Sp(n)$ is incompatible.
  \item $U(n)$ preserves the complex inner product, which is
  the natural structure of the ledger.  Restricting to
  determinant~1 gives $SU(n)$.
\end{itemize}
\end{proof}

%% ============================================================
\section{Anomaly Cancellation}
\label{sec:anomaly}
%% ============================================================

Chiral gauge theories can suffer from gauge anomalies that
spoil consistency.  The Standard Model is anomaly-free.

\begin{proposition}[Anomaly Cancellation --- Verification]
\label{thm:anomaly}
With three generations and the Standard Model charge assignments,
all gauge anomalies cancel.
\end{proposition}

\begin{proof}
The anomaly conditions per generation are:
\begin{enumerate}[label=(\roman*)]
  \item $[SU(3)]^2 U(1)$: Requires
  $\sum_{\mathrm{quarks}} Y = 0$ per generation.
  Left-handed quarks $Q_L$ ($Y = 1/6$, doublet: 2 states)
  contribute $2 \times (1/6) = 1/3$.  Right-handed up
  ($Y = 2/3$) and down ($Y = -1/3$) contribute
  $-(2/3) - (-1/3) = -1/3$.  Total: $1/3 - 1/3 = 0$.

  \item $[SU(2)]^2 U(1)$: Requires $\sum_{\mathrm{doublets}} Y = 0$.
  Quark doublet ($Y = 1/6$, $N_c = 3$ copies) and lepton
  doublet ($Y = -1/2$, 1 copy): $3 \times (1/6) + (-1/2) = 0$.

  \item $[U(1)]^3$: Requires $\sum Y^3 = 0$.  Per generation:
  $3 \times 2 \times (1/6)^3 + 3 \times (-2/3)^3 +
  3 \times (1/3)^3 + 2 \times (-1/2)^3 + (1)^3 = 0$.

  \item Gravitational-$U(1)$: Requires $\sum Y = 0$.
  Per generation: $3 \times 2 \times (1/6) + 3 \times (-2/3) +
  3 \times (1/3) + 2 \times (-1/2) + (1) = 0$.

  \item $[SU(2)]^3$: Vanishes identically because $SU(2)$
  has only pseudoreal representations.

  \item $[SU(3)]^3$: Vanishes because the chiral content
  balances: per generation, two left-handed triplets ($Q_L$
  doublet) and two right-handed triplets ($u_R$, $d_R$)
  contribute equally.
\end{enumerate}
With three generations, each generation independently satisfies
anomaly cancellation, so the total is anomaly-free.
\end{proof}

\begin{remark}[Status of the anomaly cancellation argument]
Anomaly cancellation in the SM is often cited as evidence for
the specific hypercharge assignments.  The computation above verifies
that the SM hypercharge assignments (which RS identifies as outputs of the gap
function $\mathrm{gap}(Z) = Z + Z^2 + Z^4$) satisfy anomaly cancellation.
This is a consistency check of the RS charge identification: the fact that
anomaly cancellation holds for the RS-identified charges is supporting
evidence for the identification, but not a derivation of anomaly cancellation
from RS first principles.  The connection between the gap function and the
specific hypercharge values that make anomaly cancellation automatic is a deep
structural feature of RS; a rigorous derivation is an identified open problem.
\end{remark}

%% ============================================================
\section{Two Derivations of $N_{\mathrm{gen}} = 3$}
\label{sec:two-derivations}
%% ============================================================

RS provides two independent derivations of
$N_{\mathrm{gen}} = 3$:

\begin{enumerate}
  \item \textbf{Bit decomposition}: $8 = 2^3$ tick phases
  decompose into 3 independent bits, one per spatial dimension.
  Each bit corresponds to one generation.

  \item \textbf{Face-pair count}: The 3-cube has 3 face pairs,
  and each face pair supports one generation of fermions.
\end{enumerate}

\begin{proposition}[Consistency of the Two Derivations]
\label{prop:consistency}
Both derivations give $N_{\mathrm{gen}} = D = 3$ for $D = 3$,
and they agree for \emph{all} $D$:
$\log_2(2^D) = D = F/2 = D$.  The bit count and face-pair
count are equal because the $D$-cube has $2D$ faces, giving
$D$ face pairs, and $2^D$ vertices encode $D$ bits.
\end{proposition}

\begin{remark}[Independence of the two derivations]
\label{rem:independence}
The two derivations (bit decomposition and face-pair count) are
consistent but not fully independent: both ultimately rely on the
identification $N_\mathrm{gen} = D = 3$.  The bit decomposition
uses $8 = 2^3$ (so 3 bits), and the face-pair count uses $F/2 = D = 3$;
both reduce to the single statement $D = 3$.  The two derivations offer
complementary structural perspectives, but each collapses to
``$N_\mathrm{gen} = D$'' --- which is an identification, not a derivation
of why the number of generations should equal the spatial dimension.  A
rigorous derivation of $N_\mathrm{gen} = D$ from the RS Hamiltonian
(rather than as a structural labeling) is an identified open problem.
\end{remark}

%% ============================================================
\section{Comparison with Grand Unified Theories}
\label{sec:gut}
%% ============================================================

{\small
\begin{center}
\renewcommand{\arraystretch}{1.3}
\begin{tabular}{@{}lcccp{3.2cm}@{}}
\toprule
\textbf{Framework} & \textbf{Gauge gp.} &
\textbf{$N_c$?} & \textbf{$N_{\rm gen}$?} &
\textbf{Issues} \\
\midrule
SM & Input & No & No & Postulated \\
$SU(5)$ GUT & Broken & No & No & Proton decay \\
$SO(10)$ GUT & Broken & No & No & Many parameters \\
$E_8$ & Specul. & No & No & No SM embedding \\
RS & Derived & Yes & Yes & Lie algebra open \\
\bottomrule
\end{tabular}
\end{center}
}

\begin{proposition}[No GUT Breaking Required]
\label{prop:no-gut}
All three gauge sectors originate from a single recognition
ledger.  The apparent multiplicity of gauge groups reflects
different topological aspects of the same 3-cube:
\begin{itemize}
  \item $SU(3)_C$: face-pair channels,
  \item $SU(2)_L$: spinor structure of rotations,
  \item $U(1)_Y$: overall recognition phase.
\end{itemize}
RS does not require a larger group that spontaneously breaks;
the SM group is the natural decomposition of the cube
geometry.
\end{proposition}

\begin{remark}
This implies no proton decay from GUT-scale gauge bosons (the
$X$ and $Y$ bosons of $SU(5)$), no magnetic monopoles from GUT
symmetry breaking, and no extra dimensions beyond $D = 3$.
Current proton decay bounds ($\tau_p > 2.4 \times
10^{34}$ years for $p \to e^+ \pi^0$~\cite{SuperK}) are
consistent with RS.
\end{remark}

%% ============================================================
\section{Predictions and Falsifiers}
\label{sec:predictions}
%% ============================================================

\begin{enumerate}
  \item \textbf{Exactly three colors}: $N_c = 3$, not 2 or 4.
  Any evidence of a fourth quark color would falsify RS.

  \item \textbf{Exactly three generations}: $N_{\mathrm{gen}}
  = 3$.  A fourth generation of charged fermions or light
  neutrinos would falsify RS (LEP already constrains
  $N_\nu = 2.984 \pm 0.008$~\cite{LEP}).

  \item \textbf{No additional gauge factors}: No extra
  $U(1)'$, no hidden-sector gauge group beyond what the 3-cube
  geometry provides.

  \item \textbf{No proton decay via GUT bosons}: The proton
  is stable against $SU(5)$-type decay channels.

  \item \textbf{No magnetic monopoles}: No topological defects
  from GUT symmetry breaking.
\end{enumerate}

%% ============================================================
\section{Conclusion}
\label{sec:conclusion}
%% ============================================================

The gauge group $SU(3)_C \times SU(2)_L \times U(1)_Y$ is
derived from the geometry of the $D = 3$ cube:
\begin{enumerate}
  \item $SU(3)$: from 3 face pairs of the cube (each face
  pair = one color channel).
  \item $SU(2)$: from the spinor structure of $SO(3)$
  rotations.
  \item $U(1)$: from the overall recognition phase.
  \item Unitary groups: from the complex structure of the
  ledger (8-tick forces $\mathbb{C}$).
  \item Three generations: from $D = 3$ (both the bit
  decomposition $8 = 2^3$ and the face-pair count).
  \item Anomaly cancellation: satisfied for the RS-derived
  charge assignments.
\end{enumerate}
RS is unique in deriving both the gauge group and its
representations from a single geometric structure, without
invoking grand unification or extra dimensions.

%% ============================================================
\begin{thebibliography}{99}

\bibitem{Aczel1966}
J.~Acz\'el, \emph{Lectures on Functional Equations and Their Applications},
Academic Press, New York (1966).

\bibitem{RS_Axioms}
J.~Washburn, ``The Algebra of Reality: A Recognition Science Derivation
of Physical Law,'' \emph{Axioms} \textbf{15}(2), 90 (2026).
\texttt{doi:10.3390/axioms15020090}

\bibitem{Glashow1961}
S.~L.~Glashow, ``Partial-Symmetries of Weak Interactions,''
Nucl.\ Phys.\ \textbf{22}, 579 (1961).

\bibitem{Weinberg1967}
S.~Weinberg, ``A Model of Leptons,'' Phys.\ Rev.\ Lett.\
\textbf{19}, 1264 (1967).

\bibitem{Salam1968}
A.~Salam, in \textit{Elementary Particle Theory},
ed.\ N.~Svartholm (Almqvist \& Wiksell, Stockholm, 1968),
p.~367.

\bibitem{GeorgiGlashow1974}
H.~Georgi and S.~L.~Glashow, ``Unity of All Elementary
Particle Forces,'' Phys.\ Rev.\ Lett.\ \textbf{32}, 438
(1974).

\bibitem{Georgi1975}
H.~Georgi, in \textit{Particles and Fields}, AIP Conf.\
Proc.\ \textbf{23}, 575 (1975).

\bibitem{Lisi2007}
A.~G.~Lisi, ``An Exceptionally Simple Theory of Everything,''
arXiv:0711.0770 [hep-th] (2007). [Note: preprint, not peer-reviewed;
cited as a representative $E_8$ unification proposal only.]

\bibitem{LEP}
ALEPH, DELPHI, L3, OPAL, SLD Collaborations,
``Precision Electroweak Measurements on the $Z$ Resonance,''
Phys.\ Rept.\ \textbf{427}, 257 (2006).

\bibitem{SuperK}
Super-Kamiokande Collaboration, ``Search for Proton Decay
via $p \to e^+ \pi^0$ and $p \to \mu^+ \pi^0$,''
Phys.\ Rev.\ D \textbf{95}, 012004 (2017).

\end{thebibliography}

\end{document}
